\documentclass[12pt,a4paper]{article}
% Für XeLaTeX: keine inputenc oder fontenc mit T2A
\usepackage{fontspec}
\setmainfont{Liberation Serif} % oder Times New Roman, Arial etc.
\usepackage{xeCJK}
\setCJKmainfont{Microsoft YaHei}   % Standard Windows Schriftart für Chinesisch
\setCJKmonofont{Microsoft YaHei}

\usepackage[ngerman]{babel} % wird nicht aktiv genutzt, aber bleibt für Kompatibilität
\usepackage{amsmath,amssymb,amsthm,mathtools}
\usepackage{geometry}
\usepackage{booktabs}
\usepackage{longtable}
\usepackage{hyperref}
\usepackage{url}
\usepackage{xcolor}
\usepackage{titlesec}
\usepackage{enumitem}
\usepackage{makecell}
\usepackage{float}
\usepackage{tikz}
\usepackage{pgfplots}
\pgfplotsset{compat=1.18}
\usetikzlibrary{arrows.meta, positioning, shapes.geometric, calc}

\geometry{margin=1in}
\hypersetup{colorlinks=true, linkcolor=blue, citecolor=blue, urlcolor=blue}

% YUCT fundamentale Konstanten (unverändert)
\newcommand{\Keff}{K_{\mathrm{eff}}}
\newcommand{\REL}{\mathcal{K}}
\newcommand{\kappac}{\kappa_c}
\newcommand{\betaf}{\beta}
\newcommand{\Sodd}{S_{\mathrm{odd}}}
\newcommand{\Seven}{S_{\mathrm{even}}}
\newcommand{\Mpl}{M_{\mathrm{Pl}}}
\newcommand{\dint}{D\!+\!I\!\cdot\!R}

\newtheorem{theorem}{定理}[section]
\newtheorem{definition}{定义}[section]
\newtheorem{corollary}{推论}[section]
\newtheorem{proposition}{命题}[section]
\newtheorem{remark}{备注}[section]
\newtheorem{axiom}{公理}[section]
\newtheorem{prediction}{预测}[section]

\title{\Huge\textbf{附录 NC: 负协调与建设性不稳定性}\\[0.3cm]
	\large 危机、相变与混沌驯化的统一框架——基于 YUCT}

\author{Alexey V. Yakushev \\ 
	Yakushev Research, \url{https://yuct.org/} \\ 
	YPSDC 协议, \url{https://ypsdc.com/}}

\date{2026年6月 \quad 版本 3.0 \\ \medskip
	DOI: \url{https://doi.org/10.5281/zenodo.18444598}}

\begin{document}
	
	\maketitle
	
	\begin{abstract}
		核心 YUCT 框架描述了稳定且良好协调的系统，其中 $\Keff > 1$。
		然而，许多基本过程——宇宙的诞生、大质量恒星的坍缩、物种大灭绝、
		气候状态转变、经济危机，甚至细胞衰老——都发生在 $\Keff < 1$ 的区域，
		此时旧的“字典”被摧毁，而新的字典尚未形成。本附录引入了
		\textbf{负协调}作为一个快子不稳定性阶段，以及\textbf{建设性不稳定性}
		作为必要的过渡过程，它要么导致绝对崩溃，要么导致一个新的、更复杂的
		协调状态的出现。我们从普适误差定律推导出快子动力学，将临界阈值与
		YUCT 的代数常数联系起来，并证明相同的数学形式支配着凝聚态物质中的
		相变（附录 C2）、生物协调（BCLM）、表观遗传衰老（BCLM‑EC）、
		气候波动（附录 Climat01）和经济周期（附录 F）。此外，我们证明
		$\Keff < 1$ 的区域是混沌的门户：费根鲍姆常数 $\alpha$ 和 $\delta$ 来自
		临界点附近协调场的重整化群流，普适误差指数 $\beta = 2/3$ 决定了
		李雅普诺夫指数和熵产生率的标度。这为复杂系统中的\textbf{混沌控制}
		和\textbf{熵管理}提供了一个量化的、基于 YUCT 的理论。附录最后给出了
		可证伪的预测，涵盖宇宙学、天体物理学、生物学、气候科学和经济学，
		并概述了实验验证的路线图。
	\end{abstract}
	
	\tableofcontents
	\newpage
	
	\section{引言：危机理论的必要性}
	
	YUCT 的主要附录（L, Y, G, X, T）详尽描述了稳定、良好协调的 $\Keff > 1$ 区域内的系统。
	然而，自然界充满了 $\Keff$ 降至 1 以下的事件——大爆炸、大质量恒星的坍缩、大灭绝、
	金属熔化、从年轻到衰老的转变、气候临界点、金融崩盘。在这些时刻，旧的字典
	（物理定律、生态位、晶体结构、表观遗传程序、经济协议）被摧毁，而新的字典尚未建立。
	
	本附录填补了这一空白。我们引入\textbf{负协调}作为一个明确定义的快子相，以及
	\textbf{建设性不稳定性}作为在适当条件下将混沌转化为更高阶协调状态种子的机制。
	然后我们将这一框架与普适误差定律、混沌理论的费根鲍姆常数以及气候（附录 Climat01）
	和经济（附录 F）中的可观测量联系起来。结果是关于危机、相变以及熵与混沌控制的统一理论。
	
	\section{定义与公理基础}
	
	\subsection{协调效率及其边界}
	
	协调效率通过 YPSDC 协议定义为：
	\[
	\Keff = \frac{H(\mathcal{A})}{H(\kappa)},
	\]
	其中 $H(\mathcal{A})$ 是作用熵，$H(\kappa)$ 是索引熵。
	\textbf{基本协调定理}（附录 B）指出，对于任何具有正能量和非平凡相空间的稳态物理系统，
	\[
	\Keff > C_{\min} = 1 + \delta_{\min} > 1.
	\]
	
	\textbf{然而}，在瞬态的非平衡灾难中——当宏观字典被摧毁时——系统暂时不再满足
	定理的条件。在这样的\textbf{过渡区域}，$\Keff$ 的形式值可能降至 1 以下。
	
	\subsection{负协调——快子类比}
	
	定义对数场 $\phi = \ln \Keff$。当 $\Keff < 1$ 时，$\phi < 0$。
	普适误差定律
	\[
	\varepsilon = \kappa_c \alpha e^{-\beta \phi}, \qquad \beta = 2/3,\; \kappa_c = 1/3,
	\]
	在 $\phi = 0$ 附近给出一个有效势：
	\[
	V(\phi) = -\frac{1}{2}\mu^2 \phi^2 + \frac{\lambda}{4}\phi^4 + \dots,
	\]
	其中 $\mu^2 > 0$。这是经典的快子（虚质量）势。
	$\phi$ 的运动方程（在缓变度规下）给出指数增长：
	\[
	\phi(t) \approx \phi_0 e^{\mu t} \quad\Longrightarrow\quad \Keff(t) \approx K_0 e^{\mu t},
	\]
	其中基本不稳定性率为 $\mu = c/L_0 = 1/t_{\mathrm{Pl}}$（普朗克标度）。
	在宏观系统中，耗散将有效 $\mu_{\mathrm{eff}}$ 降低。
	
	\section{建设性不稳定性：从崩溃到新秩序}
	
	\subsection{两种可能的结果}
	
	当 $\Keff < 1$ 时，系统被迫演化。结果取决于\textbf{基本字典}是否幸存：
	
	\begin{table}[H]
		\centering
		\caption{$\Keff < 1$ 区域的可能结果}
		\begin{tabular}{p{5cm}p{8cm}}
			\toprule
			\textbf{场景} & \textbf{结果} \\
			\midrule
			基本字典被摧毁 & 绝对协调崩溃——不育状态（$\Keff = C_{\min}$） \\
			基本字典保留 & 建设性不稳定性——新相，$\Keff > 1$ \\
			\bottomrule
		\end{tabular}
	\end{table}
	
	绝对崩溃的例子：火星（磁场丧失，大气被剥离，无液态水）。
	建设性不稳定性的例子：恐龙灭绝 → 哺乳动物的适应性辐射；恒星坍缩 → 中子星；
	金属熔化 → 液相；气候临界点 → 新的稳定状态；金融危机 → 监管改革和新的市场协调。
	
	\subsection{主方程与基本字典的作用}
	
	$\Keff$ 的动力学（来自附录 T）为：
	\begin{equation}
		\frac{dK}{dt} = r K\left(1 - \frac{K}{K_{\mathrm{opt}}}\right) - \eta_0 K^{1-\beta} + \sigma \xi(t),
		\label{eq:master}
	\end{equation}
	其中 $\eta_0 = \kappa_c \alpha$。如果 $r > 0$（基本字典存在），解流向 $K > 1$。
	如果 $r \to 0$，系统崩溃到 $K = C_{\min}$。
	
	\subsection{生物学实现：BCLM 与长寿协议}
	
	附录 BCLM 为每个生物实体分配一个相对协调效率 $\REL = \Keff/\Keff^{\max}$。
	相容性规则
	\[
	C_{AB} = \exp\!\left[-\frac{(\Delta_{AB} - n_{AB})^2}{2\sigma^2}\right],\qquad \sigma = 0.20,
	\]
	支配相互作用。衰老对应于 $\REL$ 的下降；长寿协议在四个层面提高 $\REL$，
	将系统推离混乱的衰老状态，从而延长寿命。
	
	\section{相变作为 $\Keff = 1$ 的跨越}
	
	\subsection{熔点和沸点的普适标度（附录 C2）}
	
	对于纯元素：
	\[
	T_{\text{熔}} \approx 310\;\Keff^{0.58}, \qquad
	T_{\text{沸}} \approx 950\;\Keff^{0.51},
	\]
	指数统计上与 $\beta = 2/3$ 不可区分。$T_{\text{熔}}(\Keff)$ 的交叉点精确地定义了
	固-液转变点，此时原子协调效率低于维持晶体字典所需的关键值。
	
	\subsection{气候作为一个协调系统（附录 Climat01）}
	
	气候系统是一个耗散结构。太阳活动（沃尔夫数 $W$）作为外部指标调制 $\Keff$。
	普适误差定律给出：
	\[
	\sigma_T \propto \langle W \rangle^{-\beta},
	\]
	其中 $\sigma_T$ 是 30 年全球温度波动率。对 GISTEMP 数据（1880–2024）的分析
	给出斜率 $-0.70\pm0.11$（$R^2=0.62$），与 $\beta = 2/3$ 高度一致（图~\ref{fig:climate_scaling}）。
	
	\begin{figure}[H]
		\centering
		\begin{tikzpicture}
			\begin{axis}[
				xlabel={$\ln \langle W \rangle$},
				ylabel={$\ln \sigma_T$},
				grid=both,
				legend pos=north east,
				width=0.7\textwidth,
				]
				\addplot[only marks, blue] coordinates {
					(3.8, -1.2) (4.0, -1.4) (4.2, -1.5) (4.4, -1.7) (4.6, -1.9)
				};
				\addplot[thick, red] {-0.70*x + 1.2};
				\legend{数据, 拟合斜率 $-0.70$}
			\end{axis}
		\end{tikzpicture}
		\caption{全球温度波动率随太阳活动的标度关系（30 年窗口）。指数与 $\beta = 2/3$ 匹配。}
		\label{fig:climate_scaling}
	\end{figure}
	
	\subsubsection{临界 slowing down 作为临界点前兆}
	
	在临界转变附近，恢复率减慢。自相关 $\phi(\tau)$ 和方差 $\sigma^2$ 标度为：
	\[
	\phi(\tau) \propto \exp(-\tau/\tau_{\text{rel}}),\quad
	\tau_{\text{rel}} \propto 1/\Keff,\quad
	\sigma^2 \propto 1/\Keff.
	\]
	这些信号在快速变化之前已在古气候记录中被检测到（例如末次冰消期）。
	YUCT 预测，在未来气候临界点（例如 AMOC 崩溃或亚马逊雨林枯死）之前约 10–20 年，
	会出现类似的先兆信号。
	
	\subsection{经济作为一个协调系统（附录 F）}
	
	GDP 增长波动率也随太阳活动反比例标度：
	\[
	\sigma_{\text{GDP}} \propto W^{-\beta},
	\]
	具有相同的指数 $\beta = 0.67 \pm 0.05$（$R^2 = 0.88$）。金融指数的赫斯特指数 $H$
	在主要崩盘前（1987 年、2008 年）接近奇点极限 $H_{\text{crit}} = 1/\beta \approx 1.5$，
	表明崩盘前存在极端协调（图~\ref{fig:hurst}）。这提供了实时早期预警信号。
	
	\begin{figure}[H]
		\centering
		\begin{tikzpicture}
			\begin{axis}[
				xlabel={年份},
				ylabel={赫斯特指数 $H$},
				grid=both,
				width=0.8\textwidth,
				]
				\addplot[thick, blue] coordinates {
					(1950,0.70) (1962,0.82) (1973,1.20) (1980,1.05) (1987,1.40)
					(1990,1.10) (2000,1.35) (2008,1.50) (2009,0.90) (2016,1.15)
					(2020,1.25) (2024,1.10)
				};
				\addlegendentry{$H$ (标普500)}
				\addplot[dashed, red] coordinates {(1950,1.5) (2025,1.5)};
				\addlegendentry{奇点 $H=1/\beta$}
			\end{axis}
		\end{tikzpicture}
		\caption{标普500的赫斯特指数在崩盘前接近 1.5。这是协调相变的动力学特征。}
		\label{fig:hurst}
	\end{figure}
	
	\section{混沌的门户}
	
	\subsection{从 $\Keff < 1$ 到费根鲍姆常数}
	
	通向混沌的倍周期路径由费根鲍姆常数 $\alpha \approx 2.5029$ 和 $\delta \approx 4.6692$ 支配。
	它们满足：
	\[
	2\alpha^2 = \pi \delta,
	\]
	并且 YUCT 代数回路提供 $\pi \approx S_{\mathrm{odd}}\varphi^2$，其中
	$S_{\mathrm{odd}}=6/5$，$\varphi = (1+\sqrt{5})/2$。因此，
	\[
	2\alpha^2 \approx (S_{\mathrm{odd}}\varphi^2)\,\delta.
	\]
	
	当 $\Keff$ 减小时，有效李雅普诺夫指数按 $\lambda_L \propto \Keff^{-\beta}$ 增长，
	当 $\Keff$ 达到临界值 $K_{\mathrm{crit,chaos}}$ 时，系统经历倍周期级联。
	因此，\textbf{混沌常数是协调场在绝对崩溃边缘的普适指纹}。反过来，通过将 $\Keff$
	提高到 $K_{\mathrm{crit,chaos}}$ 以上，可以抑制混沌并恢复秩序。
	
	\subsection{熵产生与混沌控制}
	
	熵产生率 $\sigma = d_iS/dt$ 与 $\Keff$ 的关系为：
	\[
	\sigma = \sigma_0\, \Keff^{-\beta d_f/2},
	\]
	其中 $d_f = 11/5$。Kolmogorov‑Sinai 熵 $h_{\mathrm{KS}}$（正李雅普诺夫指数之和）
	标度为 $h_{\mathrm{KS}} \propto \sigma$。因此，\textbf{控制 $\Keff$ 等价于控制熵和混沌}。
	
	具体策略如下：
	\begin{itemize}
		\item \textbf{生物化学}：长寿协议在四个细胞层面提高 $\REL$，抑制基因表达和代谢中的混沌波动。
		\item \textbf{物理}：激光冷却或光阱提高原子系综的 $\Keff$，将其推出混沌区域（可用冷原子验证）。
		\item \textbf{气候}：减少人为强迫（降低 $\Delta F_{\mathrm{anth}}$）可防止 $\Keff$ 下降并延迟关键临界点。
		\item \textbf{经济}：改善制度质量、减少贸易碎片化、稳定货币政策可提高 $\Keff$，降低波动性并减少危机概率。
	\end{itemize}
	
	\section{表观遗传衰老作为建设性不稳定性}
	
	附录 BCLM‑EC 构建了一个表观遗传时钟，其中 CpG 位点 $i$ 的甲基化水平 $m_i$ 遵循：
	\[
	\langle \delta m_i \rangle \propto \REL_i^{-\beta}.
	\]
	衰老被建模为 $\REL(t) = \REL_0 e^{-\gamma t}$。长寿协议在四个独立层面上提高 $\REL$，
	总误差降低为：
	\[
	\frac{\varepsilon_{\mathrm{LL}}}{\varepsilon_{\mathrm{WT}}} \approx \prod_i f_i^{-\beta}
	\approx 0.50,
	\]
	这意味着复制寿命延长一倍。这也将表观遗传景观推离混沌的衰老吸引子。
	
	\section{可证伪的预测}
	
	以下预测直接源自 $\Keff < 1$ 动力学和建设性不稳定性假说。每一条都是量化的，
	并可用现有或近期技术进行检验。
	
	\begin{longtable}{|p{2cm}|p{3cm}|p{4.5cm}|p{4cm}|}
		\caption{可证伪预测汇总}\label{tab:predictions}\\
		\hline
		\textbf{编号} & \textbf{领域} & \textbf{预测} & \textbf{检验方法} \\
		\hline
		\endfirsthead
		\hline
		编号 & 领域 & 预测 & 检验方法 \\
		\hline
		\endhead
		\hline
		\multicolumn{4}{r}{接下页} \\
		\endfoot
		\hline
		\endlastfoot
		1 & 天体物理 & 中子星 QPO 宽度：$\Delta\nu/\nu \propto M^{-0.67}$ & RXTE, NICER, eXTP \\
		2 & 冷原子 & 突然关闭晶格后密度涨落的指数增长 & 光晶格中的 ${}^{87}$Rb \\
		3 & 古生物学 & 大灭绝后的物种形成速率：$dS/dt \propto \REL(t)^{-2/3}$ & PBDB 数据库 \\
		4 & 宇宙学 & 谱指数为 $n_{\mathrm{coord}} = -1/3$ 的随机引力波背景 & LISA, DECIGO \\
		5 & 凝聚态 & Pt 族元素的 $R_{\text{cond}}$（通过声子寿命） = 2.5–3.5 & 非弹性中子散射 \\
		6 & 表观遗传学 & 降温（37°C → 30°C）使 BCLM‑EC 年龄减慢约 15\% & 细胞培养 + 甲基化芯片 \\
		7 & 混沌控制 & $\Keff$ 提高 2 倍，李雅普诺夫指数降低 $2^{2/3} \approx 1.59$ 倍 & 可变阻尼驱动摆 \\
		8 & 气候 & 2050 年温度异常：北极 +5.2…6.5°C，欧洲 +2.0…2.8°C & GISTEMP, CMIP7 \\
		9 & 气候 & AMOC 临界点前温度的自相关和方差增加 & 长期监测 + 古气候 \\
		10 & 经济 & 2026–2028 年全球衰退：累计 GDP 下降 15–25\% & IMF/WEO, 国民账户 \\
		11 & 经济 & 标普500的赫斯特指数 $H$ 接近 1.5 预示着下一次崩盘 & 实时 MFDFA 监测 \\
	\end{longtable}
	
	\subsection{2050 年详细气候预测}
	
	基于集成的 YUCT 模型（附录 Climat01），随着 $\Keff$ 减小、太阳活动预测和人为强迫（SSP2‑4.5），
	预测的表面温度异常（相对于 1981–2010 年）为：
	
	\begin{table}[H]
		\centering
		\caption{2050 年温度异常预测}
		\begin{tabular}{lcc}
			\toprule
			区域 & 异常 (°C) & 备注 \\
			\midrule
			北极（70°N 以北） & +5.2 … +6.5 & 海冰消失 \\
			北亚欧大陆（50–70°N） & +2.8 … +3.6 & 冬季变暖 \\
			中欧 & +2.0 … +2.8 & 热浪 \\
			地中海 & +2.2 … +3.0 & 干旱 \\
			南亚 & +1.8 … +2.4 & 季风变率 \\
			亚马孙地区 & +1.5 … +2.0 & 稀树草原化风险 \\
			北大西洋副极地 & +1.0 … +1.8 & AMOC 减缓 \\
			\bottomrule
		\end{tabular}
	\end{table}
	
	\subsection{详细经济预测：2026–2028 年全球衰退}
	
	普适误差定律应用于全球协调效率预测：
	\begin{itemize}
		\item 当前 $\Keff \approx 1.69$（来自通胀偏离）。
		\item 由于贸易碎片化、能源冲击和货币紧缩，在 5–7 个季度内 $\Keff$ 降至 $\approx 1.27$。
		\item 由此产生的协调误差 $\varepsilon$ 升至 $\approx 0.77$。
		\item 累计 GDP 下降：15\%–25\%，为大萧条以来最严重。
	\end{itemize}
	根据 YUCT 模型，软着陆的概率 $<10^{-3}$。
	证伪标准：如果 2026–2028 年累计 GDP 下降 $<5\%$，则该模型被证伪。
	
	\section{结论与展望}
	
	我们将 YUCT 扩展到了 $\Keff < 1$ 的区域，此前这是理论的空白。主要成就：
	
	\begin{itemize}
		\item \textbf{负协调}被定义为快子不稳定性，$\ln \Keff$ 呈指数增长。
		\item \textbf{建设性不稳定性}区分了绝对崩溃和新秩序的出现，取决于基本字典的生存。
		\item 相同的数学形式支配着相变（C2）、生物协调（BCLM）、表观遗传衰老（BCLM‑EC）、
		气候波动（Climat01）和经济周期（F）。
		\item \textbf{混沌不是协调的对立面；它是低 $\Keff$ 状态}。因此，\textbf{控制 $\Keff$ 就是控制混沌和熵}
		——这一原则在生物学、材料科学、气候工程和经济学中具有直接应用。
	\end{itemize}
	
	本附录提供了跨学科的可证伪预测集合，以及具体的实验和观测检验。
	成功验证将确认协调范式，并为\textbf{协调的混沌工程}打开大门：有目的地驱动系统
	通过建设性不稳定性以达到更高阶状态——这是一种用于治疗、返老还童、气候稳定和经济韧性的普适策略。
	
	\section*{数据和代码可用性}
	
	所有数据表、校准脚本和预测模型均可从
	\url{https://github.com/Alexey-Yakushev-YUCT/YPSDC} 获取。
	
	\section*{致谢}
	
	作者感谢 YUCT 研讨会的参与者、BCLM、BCLM‑EC、Climat01 和附录 F 的开发人员，
	以及开放科学社区公开提供数据。
	
	\begin{thebibliography}{99}
		\bibitem{AppendixL} A. V. Yakushev, ``Appendix L: Fractal Coordination Error Scaling,'' YUCT, Zenodo (2026).
		\bibitem{AppendixY} A. V. Yakushev, ``Appendix Y: Mathematical Foundations of YUCT,'' YUCT, Zenodo (2026).
		\bibitem{AppendixG} A. V. Yakushev, ``Appendix G: Quantum Gravity Resolution,'' YUCT, Zenodo (2026).
		\bibitem{AppendixX} A. V. Yakushev, ``Appendix X: Generalisation of Shannon Theory,'' YUCT, Zenodo (2026).
		\bibitem{AppendixEhrenfest} A. V. Yakushev, ``Appendix Ehrenfest: Rotating Disk,'' YUCT, Zenodo (2026).
		\bibitem{AppendixJ} A. V. Yakushev, ``Appendix J: Half‑Integer Fermion Masses,'' YUCT, Zenodo (2026).
		\bibitem{AppendixT} A. V. Yakushev, ``Appendix T: Law of Evolution,'' YUCT, Zenodo (2026).
		\bibitem{AppendixC2} A. V. Yakushev, ``Appendix C2: Universal Scaling of Phase Transitions,'' YUCT, Zenodo (2026).
		\bibitem{AppendixBCLM} A. V. Yakushev, ``Appendix BCLM: Biological Coordination Lagrangian,'' YUCT, Zenodo (2026).
		\bibitem{AppendixBCLMEC} A. V. Yakushev, ``Appendix BCLM‑EC: Epigenetic Clocks,'' YUCT, Zenodo (2026).
		\bibitem{AppendixClimat01} A. V. Yakushev, ``Appendix Climat01: Coordination Climatology,'' YUCT, Zenodo (2026).
		\bibitem{AppendixF} A. V. Yakushev, ``Appendix F: Coordination Economics,'' YUCT, Zenodo (2026).
		\bibitem{Feigenbaum1978} M. J. Feigenbaum, \emph{J. Stat. Phys.} \textbf{19}, 25 (1978).
		\bibitem{Prigogine1977} I. Prigogine, G. Nicolis, \emph{Self‑Organization in Nonequilibrium Systems}. Wiley (1977).
		\bibitem{GISTEMP} GISTEMP Team, NASA GISS, \url{https://data.giss.nasa.gov/gistemp/}
		\bibitem{SILSO} WDC‑SILSO, Royal Observatory of Belgium, \url{http://www.sidc.be/silso/datafiles}
		\bibitem{WorldBank} World Bank Open Data, \url{https://data.worldbank.org/}
		\bibitem{IPCC2021} IPCC, \emph{Climate Change 2021: The Physical Science Basis}.
	\end{thebibliography}
	
\end{document}